\chapter{Comparação dos Resultados do Modelo com Respostas do Público Externo}
\label{cap:resultados-extensao}

\todo{revisar}

Como parte do eixo de extensão descrito na Introdução, buscou-se comparar o desempenho do modelo desenvolvido com a capacidade humana de identificar domínios maliciosos a partir apenas da observação do nome de domínio, sem qualquer análise de tráfego ou de metadados subjacentes. Esse experimento teve como objetivo validar, empiricamente, a premissa central deste trabalho: a de que a triagem manual de URLs é insuficiente diante da sofisticação de técnicas como o \textit{typosquatting}, justificando a necessidade de um mecanismo automatizado baseado em redes neurais.

\subsection*{Metodologia da comparação}

Foi selecionado um conjunto de 200 nomes de domínio \texttt{.br}, composto por 100 domínios legítimos e 100 domínios maliciosos (relacionados a \textit{phishing} ou \textit{malware}), listados na íntegra no Apêndice~\ref{ap:formulario}. Esse conjunto foi submetido a três grupos avaliadores distintos, sob as mesmas condições:

\begin{itemize}
    \item \textbf{Público externo leigo}: participantes sem formação em Tecnologia da Informação, público-alvo da atividade de extensão, que responderam ao formulário por meio da plataforma Google Forms, totalizando 600 respostas ao longo dos 200 domínios;
    \item \textbf{Especialista}: um avaliador com formação e experiência em segurança da informação, submetido ao mesmo formulário e ao mesmo conjunto de domínios;
    \item \textbf{Modelo desenvolvido}: as versões binária e multiclasse do modelo MLP (Capítulo~\ref{cap:modelo}), que classificaram os mesmos 200 domínios a partir do tráfego DNS associado a cada um.
\end{itemize}

Em todos os casos, o enunciado apresentado foi o mesmo utilizado no formulário (Apêndice~\ref{ap:formulario}): \textit{``Ao navegar pela internet ou pelas suas redes sociais, você acessaria os links listados abaixo para ver o conteúdo?''}, com resposta binária (``Sim'' ou ``Não''). Para viabilizar a comparação direta com as métricas de classificação já empregadas ao longo deste trabalho, a resposta ``Não'' foi tratada como equivalente à classificação de um domínio como malicioso, e a resposta ``Sim'', como equivalente à classificação de um domínio como benigno.

\subsection*{Resultados obtidos}

A Tabela~\ref{tab:comparacao-publico} resume os resultados obtidos pelos três grupos avaliadores, considerando a acurácia geral e a distribuição de acertos e erros por classe: Falso Negativo (domínio malicioso julgado como ``Sim''), Falso Positivo (domínio benigno julgado como ``Não''), Verdadeiro Positivo (domínio malicioso corretamente julgado como ``Não'') e Verdadeiro Negativo (domínio benigno corretamente julgado como ``Sim'').

\begin{table}[h]
\centering
\caption{Comparação de desempenho entre avaliadores humanos e o modelo desenvolvido}
\label{tab:comparacao-publico}
\footnotesize
\resizebox{\textwidth}{!}{%
\begin{tabular}{lcccccc}
\toprule
Avaliador / Modelo & Acurácia Geral & FN & FP & VP & VN & Tempo Total \\
\midrule
Especialista            & 36,50\% (73/200)  & 67,00\% (67) & 60,00\% (60) & 33,00\% (33) & 40,00\% (40) & 8 min 24 s \\
Média do público externo & 43,00\%          & 57,75\%      & 56,25\%      & 42,25\%      & 43,75\%      & 11 min 36 s \\
Modelo Binário           & 93,00\% (186/200) & 2,00\% (2)   & 12,00\% (12) & 98,00\% (98) & 88,00\% (88) & 9,77 s \\
Modelo Multiclasse       & 95,00\% (190/200) & 4,00\% (4)   & 6,00\% (6)   & 96,00\% (96) & 94,00\% (94) & 9,41 s \\
\bottomrule
\end{tabular}%
}
\end{table}

\subsection*{Análise dos resultados}

O resultado mais evidente é a diferença de magnitude entre a acurácia dos avaliadores humanos e a do modelo desenvolvido: enquanto o modelo binário e o modelo multiclasse classificaram corretamente 93\% e 95\% dos 200 domínios, respectivamente, o especialista acertou apenas 36,50\% dos casos, e a média do público externo, 43,00\%. Como o conjunto de teste era balanceado entre domínios legítimos e maliciosos, uma classificação aleatória resultaria, em expectativa, em 50\% de acerto -- ou seja, tanto o especialista quanto o público em geral tiveram desempenho pior do que o esperado para uma escolha puramente aleatória. Isso indica que o erro humano observado não decorre apenas de incerteza, mas de um viés sistemático: domínios maliciosos com aparência cuidadosamente forjada tendem a inspirar mais confiança do que muitos domínios legítimos, porém pouco conhecidos.

Esse viés fica particularmente evidente nas taxas de Falso Negativo, isto é, domínios maliciosos julgados como seguros: 67,00\% para o especialista e 57,75\% em média para o público externo. Em outras palavras, a maioria dos domínios maliciosos apresentados no formulário enganou tanto o avaliador leigo quanto o avaliador com formação em segurança da informação, o que reforça a hipótese discutida na Introdução deste trabalho de que o \textit{typosquatting} e técnicas correlatas são projetados especificamente para explorar os padrões de confiança usados por humanos ao avaliar a legitimidade visual de um nome de domínio. Ao mesmo tempo, as taxas de Falso Positivo (60,00\% para o especialista e 56,25\% para o público) mostram que essa dificuldade não se restringe a subestimar o risco: uma parcela expressiva de domínios legítimos, porém com nomes incomuns ou pouco familiares, também foi julgada como suspeita, sugerindo que o julgamento humano nesse cenário se aproxima mais de uma resposta enviesada por familiaridade do que de uma avaliação técnica consistente.

O modelo desenvolvido, por outro lado, manteve taxas de Falso Negativo muito baixas (2,00\% no cenário binário e 4,00\% no multiclasse), o que indica que o mecanismo proposto captura a esmagadora maioria dos domínios maliciosos que conseguiram enganar os avaliadores humanos. A comparação entre as duas versões do modelo também revela uma pequena diferença de perfil: o modelo binário apresentou Falso Negativo ligeiramente menor (2,00\% contra 4,00\%), e, consequentemente, Verdadeiro Positivo mais alto (98,00\% contra 96,00\%), enquanto o modelo multiclasse apresentou Falso Positivo menor (6,00\% contra 12,00\%) e Verdadeiro Negativo mais alto (94,00\% contra 88,00\%), sugerindo maior precisão ao reconhecer domínios legítimos, ao custo de uma fração um pouco maior de domínios maliciosos não detectados.

Por fim, a comparação dos tempos totais de avaliação reforça a inviabilidade prática da triagem manual em escala: o especialista levou 8 minutos e 24 segundos para avaliar os 200 domínios (aproximadamente 2,52 segundos por domínio), e o público externo, em média, 11 minutos e 36 segundos (aproximadamente 3,48 segundos por domínio). O modelo, por sua vez, classificou o mesmo conjunto em menos de 10 segundos no total (aproximadamente 48 milissegundos por domínio na versão binária e 47 milissegundos na multiclasse) -- uma diferença de mais de 50 vezes em relação ao especialista e de mais de 70 vezes em relação ao público externo. Essa diferença de velocidade é tão relevante quanto a diferença de acurácia: um sistema de detecção de intrusão precisa operar no momento em que uma consulta DNS é realizada, o que é, por definição, incompatível com qualquer forma de revisão manual.

Em conjunto, esses resultados corroboram, por meio de evidência empírica com participação humana, a premissa que motivou este trabalho: a avaliação da legitimidade de um domínio a partir de sua aparência superficial é uma tarefa na qual o julgamento humano se mostra pouco confiável e lento demais para ser aplicado em escala, o que evidencia a necessidade de um mecanismo automatizado, baseado em características comportamentais do tráfego DNS, como o proposto neste trabalho.